\documentclass[11pt]{article}
\usepackage[a4paper,margin=1in]{geometry}
\usepackage[T1]{fontenc}
\usepackage{lmodern}
\usepackage{amsmath}
\usepackage{booktabs}
\usepackage{longtable}
\usepackage[numbers,sort&compress]{natbib}
\usepackage[colorlinks=true,linkcolor=blue,citecolor=blue,urlcolor=blue]{hyperref}

\title{Supplementary Information\\
\large Cognate redocking as a quality-control step for cryo-EM kinase
receptor structures: a study of ATR/CHK1 and CDK-activating kinase}
\author{}
\date{}

\begin{document}
\maketitle

\section*{Supplementary Section S1: Other CASF-2016 evaluations report
higher top-pose success rates under a different protocol}

The 58\% figure cited in the main text (Methods, Results \S3.3) is
AutoDock Vina's \emph{self-docking} (redocking) success rate on the
285-complex PDBbind Core Set, as reported by R\"ohrig et
al.~\cite{zhang2023}: the ligand is removed from its own co-crystallized
receptor and re-docked from scratch, then scored against a $\le$2.0~\AA{}
RMSD threshold to the native pose -- the same protocol used throughout
this study. This is a different, and considerably harder, test than
CASF-2016's own \emph{docking power} metric, which asks whether the
correct (near-native) pose is ranked at or near the top of a
pre-generated, decoy-enriched pose ensemble rather than found by a full
conformational search. Under that docking-power protocol, Su et
al.~\cite{su2019} report that AutoDock Vina achieves a top-1 pose success
rate of 90.2\% (257 of 285 complexes) on CASF-2016. The two figures (58\%
self-docking vs.\ 90.2\% docking power) are not directly comparable --
they test different things -- but the gap illustrates why the 58\%
self-docking baseline, not the 90.2\% docking-power figure, is the
appropriate comparator for this study's redocking protocol, and why the
58\% value should not be read as a general ``Vina success rate.''

\section*{Supplementary Table S2: Residue-deletion list (flexible-sidechain
receptor preparation)}

Methods \S2.1 and \S2.8 report that Meeko's receptor-preparation step
automatically deleted approximately 400 residues from the 9L40 receptor
whose side-chain density could not be matched against a template at the
structure's resolution, and that all deleted residues were $>$10~\AA{}
from the docking box. \textbf{This enumerated per-residue list is not
currently available as a standalone data file in the project
repository}: the deletions are performed programmatically by Meeko at
receptor-preparation time and were not captured to a separate log or
table when the pipeline was originally run. Reproducing this list exactly
would require re-running \texttt{mk\_prepare\_receptor.py} on the 9L40
structure and diffing its output against the input receptor -- the
authors have not done this and are flagging the gap here rather than
presenting a reconstructed or estimated list as original data. The
prepared receptor and ligand PDBQT files used for the flexible-docking
runs (from which this list could be regenerated) are available in the
pipeline repository (Data availability, main text).

\section*{Supplementary Table S4: Complete developability-passed
derivative ranking}

The complete ranking (677 developability-passed derivatives, satisfying $\ge$3/4 drug-likeness rules with no structural alert) is provided as \texttt{Table\_S4\_derivative\_ranking.csv} alongside this document. The first 20 ranked entries are shown below for reference; compounds flagged \texttt{selected\_for\_pilot\_docking = yes} are the 35 derivatives that entered the 40-compound pilot library (Table 4, main text).

\emph{Note on the main-text count:} the Methods section of the main text states 680 of 786 candidates satisfied $\ge$3/4 rules; the current pipeline output (\texttt{developability\_predictions.csv}) shows 677. This small discrepancy has not been resolved and is reported here rather than silently corrected in the main text.

\begin{center}
\footnotesize
\begin{tabular}{rlrrr}
\toprule
Rank & Compound ID (if selected) & Mol.\ wt.\ (Da) & Rules passed & Selected \\
\midrule
1 & deriv\_000 & 331.4 & 4/4 & yes \\
2 & deriv\_001 & 332.4 & 4/4 & yes \\
3 & deriv\_002 & 340.4 & 4/4 & yes \\
4 & deriv\_003 & 341.4 & 4/4 & yes \\
5 & deriv\_004 & 345.4 & 4/4 & yes \\
6 & deriv\_005 & 345.4 & 4/4 & yes \\
7 & deriv\_006 & 346.4 & 4/4 & yes \\
8 & deriv\_007 & 348.4 & 4/4 & yes \\
9 & deriv\_008 & 348.4 & 4/4 & yes \\
10 & deriv\_009 & 349.4 & 4/4 & yes \\
11 & deriv\_010 & 349.4 & 4/4 & yes \\
12 & deriv\_011 & 349.4 & 4/4 & yes \\
13 & deriv\_012 & 349.4 & 4/4 & yes \\
14 & deriv\_013 & 350.4 & 4/4 & yes \\
15 & deriv\_014 & 350.4 & 4/4 & yes \\
16 & deriv\_015 & 350.4 & 4/4 & yes \\
17 & deriv\_016 & 350.4 & 4/4 & yes \\
18 & deriv\_017 & 351.4 & 4/4 & yes \\
19 & deriv\_018 & 351.4 & 4/4 & yes \\
20 & deriv\_019 & 354.4 & 4/4 & yes \\
\bottomrule
\end{tabular}
\end{center}


\section*{Supplementary Table S5: Flexible-sidechain redocking, individual
replicate values}

Individual-seed naive and symmetry-corrected RMSDs and Vina scores for
the flexible-sidechain redocking run (five residues, 30~\AA{} box, Meeko
preparation; Methods \S2.8, Results \S3.4). The matched rigid control
(same 30~\AA{} box and Meeko-prepared receptor/ligand, no residues
declared flexible) is reported in the main text and pipeline repository
as a mean $\pm$ SD (4.50 $\pm$ 0.05~\AA{}); its individual per-seed values
were not retained as a separate structured output and are not
reproduced here for the same reason given for Table S2.

\begin{center}
\begin{tabular}{lrrrr}
\toprule
Condition & Seed & Naive RMSD (\AA) & Symm.\ RMSD (\AA) & Vina score (kcal/mol) \\
\midrule
Flexible (5 res., 30~\AA{} box) & 1000 & 4.5427 & 4.2882 & -9.795 \\
Flexible (5 res., 30~\AA{} box) & 1001 & 2.7159 & 2.2256 & -9.547 \\
Flexible (5 res., 30~\AA{} box) & 1002 & 2.774 & 2.4945 & -9.233 \\
Flexible (5 res., 30~\AA{} box) & mean+/-SD &  & 3.0028 +/- 0.9155 &  \\
\bottomrule
\end{tabular}

\end{center}

\section*{Supplementary Table S6: Individual replicate RMSDs underlying
Tables 1--3}

Symmetry-corrected RMSD for every seeded replicate (three seeds per
structure) behind the summary values reported in Tables 1--3 of the main
text: the three primary ATR/CHK1 structures (exhaustiveness 32), the 14
additional CAK survey structures, and the 20 CASF-2016 sanity-check
complexes (all at exhaustiveness 16, per Methods).

{\footnotesize
\begin{longtable}{llrr}
\toprule
Table & Structure & Seed & Symm.\ RMSD (\AA) \\
\midrule
\endhead
Table 1 (primary, exhaustiveness 32) & 9L40 & 1000 & 2.5334 \\
Table 1 (primary, exhaustiveness 32) & 9L40 & 1001 & 2.5183 \\
Table 1 (primary, exhaustiveness 32) & 9L40 & 1002 & 2.541 \\
Table 1 (primary, exhaustiveness 32) & 9L4B & 1000 & 0.586 \\
Table 1 (primary, exhaustiveness 32) & 9L4B & 1001 & 0.6217 \\
Table 1 (primary, exhaustiveness 32) & 9L4B & 1002 & 0.6242 \\
Table 1 (primary, exhaustiveness 32) & 2YM8 & 1000 & 0.7503 \\
Table 1 (primary, exhaustiveness 32) & 2YM8 & 1001 & 0.7596 \\
Table 1 (primary, exhaustiveness 32) & 2YM8 & 1002 & 0.7737 \\
Table 2 (17-structure survey) & 8P6V & 1000 & 6.1362 \\
Table 2 (17-structure survey) & 8P6V & 1001 & 6.1415 \\
Table 2 (17-structure survey) & 8P6V & 1002 & 6.1382 \\
Table 2 (17-structure survey) & 8P6W & 1000 & 7.3217 \\
Table 2 (17-structure survey) & 8P6W & 1001 & 7.3382 \\
Table 2 (17-structure survey) & 8P6W & 1002 & 7.3151 \\
Table 2 (17-structure survey) & 8P6X & 1000 & 6.5755 \\
Table 2 (17-structure survey) & 8P6X & 1001 & 6.6673 \\
Table 2 (17-structure survey) & 8P6X & 1002 & 6.5535 \\
Table 2 (17-structure survey) & 8P6Z & 1000 & 2.1761 \\
Table 2 (17-structure survey) & 8P6Z & 1001 & 6.1707 \\
Table 2 (17-structure survey) & 8P6Z & 1002 & 6.1794 \\
Table 2 (17-structure survey) & 8P70 & 1000 & 6.2449 \\
Table 2 (17-structure survey) & 8P70 & 1001 & 6.2502 \\
Table 2 (17-structure survey) & 8P70 & 1002 & 6.2477 \\
Table 2 (17-structure survey) & 8P71 & 1000 & 7.0164 \\
Table 2 (17-structure survey) & 8P71 & 1001 & 7.0119 \\
Table 2 (17-structure survey) & 8P71 & 1002 & 7.0354 \\
Table 2 (17-structure survey) & 8P72 & 1000 & 2.6997 \\
Table 2 (17-structure survey) & 8P72 & 1001 & 2.6968 \\
Table 2 (17-structure survey) & 8P72 & 1002 & 2.7035 \\
Table 2 (17-structure survey) & 8P73 & 1000 & 5.8317 \\
Table 2 (17-structure survey) & 8P73 & 1001 & 5.842 \\
Table 2 (17-structure survey) & 8P73 & 1002 & 5.8357 \\
Table 2 (17-structure survey) & 8P74 & 1000 & 6.6669 \\
Table 2 (17-structure survey) & 8P74 & 1001 & 6.675 \\
Table 2 (17-structure survey) & 8P74 & 1002 & 6.6784 \\
Table 2 (17-structure survey) & 8P75 & 1000 & 1.0119 \\
Table 2 (17-structure survey) & 8P75 & 1001 & 1.0533 \\
Table 2 (17-structure survey) & 8P75 & 1002 & 1.0206 \\
Table 2 (17-structure survey) & 8P76 & 1000 & 6.6039 \\
Table 2 (17-structure survey) & 8P76 & 1001 & 6.5926 \\
Table 2 (17-structure survey) & 8P76 & 1002 & 6.5984 \\
Table 2 (17-structure survey) & 8P77 & 1000 & 3.0258 \\
Table 2 (17-structure survey) & 8P77 & 1001 & 6.0337 \\
Table 2 (17-structure survey) & 8P77 & 1002 & 6.0076 \\
Table 2 (17-structure survey) & 8P78 & 1000 & 2.9874 \\
Table 2 (17-structure survey) & 8P78 & 1001 & 2.9735 \\
Table 2 (17-structure survey) & 8P78 & 1002 & 2.9702 \\
Table 2 (17-structure survey) & 8P7L & 1000 & 6.1469 \\
Table 2 (17-structure survey) & 8P7L & 1001 & 6.1418 \\
Table 2 (17-structure survey) & 8P7L & 1002 & 6.199 \\
Table 3 (CASF-2016 sanity check) & 1O3F & 1000 & 1.4579 \\
Table 3 (CASF-2016 sanity check) & 1O3F & 1001 & 1.4884 \\
Table 3 (CASF-2016 sanity check) & 1O3F & 1002 & 1.4151 \\
Table 3 (CASF-2016 sanity check) & 1SYI & 1000 & 11.8023 \\
Table 3 (CASF-2016 sanity check) & 1SYI & 1001 & 11.8423 \\
Table 3 (CASF-2016 sanity check) & 1SYI & 1002 & 11.8268 \\
Table 3 (CASF-2016 sanity check) & 2CET & 1000 & 7.7526 \\
Table 3 (CASF-2016 sanity check) & 2CET & 1001 & 7.7367 \\
Table 3 (CASF-2016 sanity check) & 2CET & 1002 & 7.742 \\
Table 3 (CASF-2016 sanity check) & 2J7H & 1000 & 9.3411 \\
Table 3 (CASF-2016 sanity check) & 2J7H & 1001 & 9.3751 \\
Table 3 (CASF-2016 sanity check) & 2J7H & 1002 & 9.3633 \\
Table 3 (CASF-2016 sanity check) & 2P4Y & 1000 & 0.5712 \\
Table 3 (CASF-2016 sanity check) & 2P4Y & 1001 & 0.4968 \\
Table 3 (CASF-2016 sanity check) & 2P4Y & 1002 & 0.4805 \\
Table 3 (CASF-2016 sanity check) & 2W4X & 1000 & 1.6605 \\
Table 3 (CASF-2016 sanity check) & 2W4X & 1001 & 1.6689 \\
Table 3 (CASF-2016 sanity check) & 2W4X & 1002 & 1.6799 \\
Table 3 (CASF-2016 sanity check) & 2XBV & 1000 & 1.0321 \\
Table 3 (CASF-2016 sanity check) & 2XBV & 1001 & 1.0494 \\
Table 3 (CASF-2016 sanity check) & 2XBV & 1002 & 1.0399 \\
Table 3 (CASF-2016 sanity check) & 3ARQ & 1000 & 3.1997 \\
Table 3 (CASF-2016 sanity check) & 3ARQ & 1001 & 3.2081 \\
Table 3 (CASF-2016 sanity check) & 3ARQ & 1002 & 3.1914 \\
Table 3 (CASF-2016 sanity check) & 3B65 & 1000 & 0.486 \\
Table 3 (CASF-2016 sanity check) & 3B65 & 1001 & 0.4871 \\
Table 3 (CASF-2016 sanity check) & 3B65 & 1002 & 0.4886 \\
Table 3 (CASF-2016 sanity check) & 3D6Q & 1000 & 6.1551 \\
Table 3 (CASF-2016 sanity check) & 3D6Q & 1001 & 6.1653 \\
Table 3 (CASF-2016 sanity check) & 3D6Q & 1002 & 6.1546 \\
Table 3 (CASF-2016 sanity check) & 3DX1 & 1000 & 3.3101 \\
Table 3 (CASF-2016 sanity check) & 3DX1 & 1001 & 3.2738 \\
Table 3 (CASF-2016 sanity check) & 3DX1 & 1002 & 3.294 \\
Table 3 (CASF-2016 sanity check) & 3KWA & 1000 & 3.3611 \\
Table 3 (CASF-2016 sanity check) & 3KWA & 1001 & 3.3752 \\
Table 3 (CASF-2016 sanity check) & 3KWA & 1002 & 3.5037 \\
Table 3 (CASF-2016 sanity check) & 3NQ9 & 1000 & 5.6598 \\
Table 3 (CASF-2016 sanity check) & 3NQ9 & 1001 & 9.8728 \\
Table 3 (CASF-2016 sanity check) & 3NQ9 & 1002 & 6.0619 \\
Table 3 (CASF-2016 sanity check) & 3PWW & 1000 & 1.6869 \\
Table 3 (CASF-2016 sanity check) & 3PWW & 1001 & 1.7456 \\
Table 3 (CASF-2016 sanity check) & 3PWW & 1002 & 1.7894 \\
Table 3 (CASF-2016 sanity check) & 3WTJ & 1000 & 6.3503 \\
Table 3 (CASF-2016 sanity check) & 3WTJ & 1001 & 6.3478 \\
Table 3 (CASF-2016 sanity check) & 3WTJ & 1002 & 6.3491 \\
Table 3 (CASF-2016 sanity check) & 4IH7 & 1000 & 4.5736 \\
Table 3 (CASF-2016 sanity check) & 4IH7 & 1001 & 4.572 \\
Table 3 (CASF-2016 sanity check) & 4IH7 & 1002 & 4.5719 \\
Table 3 (CASF-2016 sanity check) & 4K18 & 1000 & 4.5222 \\
Table 3 (CASF-2016 sanity check) & 4K18 & 1001 & 4.4707 \\
Table 3 (CASF-2016 sanity check) & 4K18 & 1002 & 4.5145 \\
Table 3 (CASF-2016 sanity check) & 4QAC & 1000 & 0.5252 \\
Table 3 (CASF-2016 sanity check) & 4QAC & 1001 & 0.5617 \\
Table 3 (CASF-2016 sanity check) & 4QAC & 1002 & 0.5207 \\
Table 3 (CASF-2016 sanity check) & 4W9C & 1000 & 9.5718 \\
Table 3 (CASF-2016 sanity check) & 4W9C & 1001 & 9.6125 \\
Table 3 (CASF-2016 sanity check) & 4W9C & 1002 & 9.5754 \\
Table 3 (CASF-2016 sanity check) & 4W9I & 1000 & 1.2615 \\
Table 3 (CASF-2016 sanity check) & 4W9I & 1001 & 1.4075 \\
Table 3 (CASF-2016 sanity check) & 4W9I & 1002 & 1.4124 \\
\bottomrule
\end{longtable}

}

\bibliographystyle{unsrtnat}
\bibliography{refs}

\end{document}
